\documentclass[11pt, a4paper, oneside]{article}
\usepackage[ngerman]{babel}
\usepackage{worksheet}

\begin{document}
	\author{L. Bung}
	\title{Assoziation, Aggregation \hspace{10cm} und Komposition}
	\subject{SAE}
	\class{E2FI}
	\maketitle
	
	\begin{hint}{Verknüpfungen zwischen Objekten}
		In vielen Fällen ist es nötig, dass Objekte Informationen über andere Objekte benötigen oder auf diese verweisen.
		Beispiel: Einem Mitarbeiter ist ein bestimmter Dienstwagen zugeordnet.\\
		
		Die Verbindung zwischen den beiden Klassen nennt man eine \textbf{Assoziation}.
		In UML werden Assoziationen durch einen Pfeil mit offener Spitze dargestellt.
		Die Assoziation kann in eine Richtung (\emph{unidirektional}) oder in beide Richtungen (\emph{bidirektional}) gerichtet sein -- je nachdem werden auch die Pfeilspitzen gesetzt.
		Außerdem können -- ähnlich zu ER-Modellen -- Kardinalitäten angegeben werden.
		
		\begin{center}
			\begin{tikzpicture}
				\begin{class}[text width=4cm]{Fahrzeug}{7,0}
					\attribute{-kennzeichen: str}
					\attribute{-leistung: int}
					\operation{+tanken(liter: int)}
				\end{class}
				\begin{class}[text width=4cm]{Mitarbeiter}{0,0}
					\attribute{-name: str}
					\attribute{-dienstwagen: Fahrzeug}
					\operation{+arbeiten(stunden: int)}
				\end{class}
				\unidirectionalAssociation{Mitarbeiter}{0..1}{}{Fahrzeug}
			\end{tikzpicture}
		\end{center}
		
		Eine stärkere Verbindung ist die \textbf{Aggregation}.
		Hier sind zwei Objekte Teil eines Ganzen -- zum Beispiel besteht eine Ehe aus zwei Ehepartnern.
		Wichtig: Die Ehepartner bestehen auch nach Auflösung der Ehe weiter, existieren also unabhängig von ihr.
		Die Aggregation wird in UML durch eine leere Raute dargestellt.
		
		\begin{center}
			\begin{tikzpicture}
				% Source - https://tex.stackexchange.com/a/391476
				% Posted by Torbjørn T.
				% Retrieved 2026-05-19, License - CC BY-SA 3.0
				
				\renewcommand{\aggregation}[4]
				{
					\draw[umlcd style, -open diamond] (#4) -- (#1)
					node[near end, above]{#2}
					node[near end, below]{#3};
				}
				
				\begin{class}[text width=4cm]{Person}{0,0}
					\attribute{-vorname: str}
					\attribute{-nachname: str}
				\end{class}
				\begin{class}[text width=4cm]{Ehe}{7,0}
					\attribute{+partner1: Person}
					\attribute{+partner2: Person}
				\end{class}
				\aggregation{Ehe}{}{}{Person}
			\end{tikzpicture}
		\end{center}
		
		Noch stärker ist die \textbf{Komposition}.
		Der Unterschied zur Aggregation liegt darin, dass die Teilobjekte nicht ohne das große Ganze existieren können.
		Ein Beispiel hierfür wäre ein Gebäude, das aus mehreren Räumen besteht -- jedoch existieren die Räume auch nicht mehr, wenn das Gebäude abgerissen wird.
		Die Komposition wird in UML durch eine ausgefüllte Raute dargestellt.
		
		\begin{center}
			\begin{tikzpicture}
				% Source - https://tex.stackexchange.com/a/391476
				% Posted by Torbjørn T.
				% Retrieved 2026-05-19, License - CC BY-SA 3.0
				
				\renewcommand{\composition}[4]
				{
					\draw[umlcd style, fill=\umldrawcolor, -diamond] (#4) -- (#1)
					node[near end, above]{#2}
					node[near end, below]{#3};
				}
				
				\begin{class}[text width=4cm]{Raum}{0,0}
					\attribute{-nummer: int}
					\operation{+aufschließen()}
				\end{class}
				\begin{class}[text width=4cm]{Gebäude}{7,0}
					\attribute{+adresse: str}
					\attribute{-räume: list(Raum)}	
				\end{class}
				\composition{Gebäude}{}{}{Raum}
			\end{tikzpicture}
		\end{center}
	\end{hint}
	
	\pagebreak
	\task{Assoziation, Aggregation oder Komposition?}
	
	a) Entscheiden Sie jeweils, was für eine Art von Verbindung vorliegt.
	
	\begin{table}[h!]
		\begin{tabularx}{\textwidth}{|l|X|l|}
			\hline
			\textbf{Beispiel} & \textbf{Art der Verbindung} & \textbf{Kardinalität}\\
			\hline
			\hline
			Bibliothek $\longleftrightarrow$ Bücher &&\\
			\hline
			Student $\longleftrightarrow$ Kurs &&\\
			\hline
			Computer $\longleftrightarrow$ CPU &&\\
			\hline
			Fahrrad $\longleftrightarrow$ Räder &&\\
			\hline
			Playlist $\longleftrightarrow$ Songs &&\\
			\hline
			Flughafen $\longleftrightarrow$ Flugzeuge &&\\
			\hline
			Betriebssystem $\longleftrightarrow$ Prozesse &&\\
			\hline
			Git-Repository $\longleftrightarrow$ Commits &&\\
			\hline
		\end{tabularx}
	\end{table}
	
	b) Finden Sie selbst weitere Beispiele für Assoziationen, Aggregationen und Kompositionen aus verschiedenen Bereichen.
	Überlegen Sie sich auch, welche Kardinalität die Verbindung jeweils hat.
	
	\task{Umsetzung in Python}
	
	\begin{enumerate}[label=\alph*)]
		\item Recherchieren Sie, wie man Assoziationen, Aggregationen und Kompositionen in Python umsetzen kann.
		Schreiben Sie für jeden Fall ein Beispiel-Programm.
		Erklären Sie in Ihrem Programm durch Kommentare, wie das Konzept umgesetzt wird.
		
		\item Setzen Sie das folgende UML-Diagramm in Python um:
		\begin{figure}[h!]
			\centering
			\begin{tikzpicture}
				% Source - https://tex.stackexchange.com/a/391476
				% Posted by Torbjørn T.
				% Retrieved 2026-05-19, License - CC BY-SA 3.0
				
				\renewcommand{\composition}[4]
				{
					\draw[umlcd style, fill=\umldrawcolor, -diamond] (#4) -- (#1)
					node[near end, above]{#2}
					node[near end, below]{#3};
				}
				\begin{class}[text width=5cm]{Bestellposition}{0,-4}
					\attribute{-produkt: str}
					\attribute{-menge: int}
					\operation{-\_\_str\_\_(): str}
				\end{class}
				\begin{class}[text width=7cm]{Bestellung}{0,0}
					\attribute{-positionen: list(Bestellposition)}
					\operation{+add\_position(produkt: str, menge: int)}
					\operation{+get\_positionen(): str}
				\end{class}
				\composition{Bestellung}{}{}{Bestellposition}
			\end{tikzpicture}
		\end{figure}
	\end{enumerate}
	
\end{document}